\documentclass[10pt]{article}
\usepackage[margin=20mm]{geometry}
\usepackage{amsmath,amssymb,fontspec,unicode-math,xurl,hyperref}
\IfFontExistsTF{DejaVu Serif}{\setmainfont{DejaVu Serif}}{\setmainfont{Latin Modern Roman}}
\setmathfont{latinmodern-math.otf}
\setlength{\parindent}{0pt}\setlength{\parskip}{6pt}\sloppy
\title{Technical review of high-dimensional sequential prediction}
\author{AI-assisted technical review}
\date{}
\begin{document}
\maketitle
Abstract: A closed-convex one-event counterexample invalidates the accelerated appendix reduction when events do not cover the prediction space. Main-text expectation, comparator and logarithmic-factor steps also need correction. Explicit proof repairs preserve the principal regret orders up to logarithms under the intended information timing, while exact rates and the general accelerated oracle guarantee remain unresolved.\par
\section*{Identity, scope and outcome}
Georgy Noarov, Ramya Ramalingam, Aaron Roth and Stephan Xie, High-Dimensional Prediction for Sequential Decision Making, ICML 2025, PMLR 267:46762–46783, OpenReview `uRAgIVnAO6`. The reviewed artifact is the frozen 22-page publisher PDF. Main pp. 1–9 and Appendices A–D, pp. 13–22, were read. The decisive formulas and algorithms were inspected in rendered physical pages. The paper is theoretical: there are no experimental figures, tables or reported training runs to replicate. The supplied bounded search did not identify author code; this is a scoped non-discovery.\par
The independently checked assessment is Major for the invalid core reductions in P1 and P2, with additional Minor corrections and evidence requests about exact rates and computational assumptions. A counterexample disproves the Appendix D candidate-set reduction under its stated disjointness assumption. That accelerated variant is separate from the general inner-loop algorithm. Main-text proofs also contain an invalid expectation identity and missing factors, but their principal bounds up to logarithms admit the explicit repairs below. Repairability and survival of other results do not validate those core derivations or lower their severity. These checks do not establish that all qualitative high-dimensional prediction and decision-making conclusions are false, nor do they certify every claimed optimal rate or priority claim.\par
\section*{P1: Disjoint events need not cover the prediction space}
Appendix D, p. 19, assumes binary events with \(\sum_E E(x,\widehat s)\leq1\). It retains one minimizer \(\widehat s^{*,E}\) from each active region and claims in Lemma D.4 that restricting to these points preserves the minimax problem. This omits the case in which no event is active.\par
Take \(S=[0,1]\), one event \(E(\widehat s)=\mathbf1[\widehat s\geq1/2]\), and positive weights \(q_+=3/4\), \(q_-=1/4\). The event is binary and disjoint, and its active region \([1/2,1]\) is nonempty, closed and convex, with an elementary separation oracle. The surrogate is \[u(\widehat s,s)=\tfrac12\mathbf1[\widehat s\geq1/2](\widehat s-s).\] Definition D.1 returns the unique active-region minimizer \(\widehat s^{*,E}=1/2\), so Algorithm 3 is restricted to \(P=\{1/2\}\). Its minimax value is \(\max_{s\in[0,1]}u(1/2,s)=1/4\). In the full problem, the point prediction \(\widehat s=0\) has payoff zero for every \(s\). Conversely \(s=0\) makes every payoff nonnegative, proving that the full minimax value is exactly zero. The gap is \(1/4\); a solution supported only on \(P\) cannot be \(\epsilon\)-optimal for \(\epsilon<1/4\). Any positive imbalance \(q_+>q_-\) gives the same obstruction, scaled accordingly.\par
An independently written linear program reproduces values \(0\), \(0.25\) and \(0\) for the full candidate set, the printed restricted set, and the set repaired by adding the null point. The analytic argument proves the continuum result, so the conclusion does not depend on discretization. Lemma D.4’s proposed mass transfer on p. 22 maps the full optimum \(\delta_0\) to total mass zero, exposing the precise failure: it is not a probability distribution. Its printed dominance inequality also has the minimization direction reversed.\par
One repair is to assume a genuine exhaustive partition. Another is to include an accessible representative of the region on which all original events vanish; its original surrogate payoff is identically zero. Finding such a representative and handling empty event regions require a stated oracle or additional geometric assumptions. Merely adding a new positively weighted event would change the game. Requiring a cover also narrows the theorem’s stated scope.\par
Remark D.2 needs a separate boundary qualification: lexicographic best-response cells need not be closed polytopes. On \([0,1]\), utilities \(u_1(s)=0\), \(u_2(s)=s\), with ties assigned to action 1, give the action-2 cell \((0,1]\). Minimizing a positive linear function on this cell has an unattained infimum. Closed-cell optimization, approximate feasible representatives and tie-breaking must be reconciled before claiming exact solutions to Definition D.1. The general Algorithm 1 allows discontinuous events and does not require this Appendix D reduction; therefore the counterexample is localized to the accelerated method, rather than a refutation of the main prediction algorithm. The broader repair and its bit-complexity guarantee remain unverified. P1 is Major because the accelerated contribution relies on a false minimax equivalence and an algorithm that fails its stated approximation guarantee; localization to an appendix does not reduce that defect to a minor one.\par
\section*{P2: Repair the expectation and comparator steps in the decision reductions}
Definition 2.2 defines \(\mathrm{Bias}_T(E)=\max_i\mathbb E|Z_{E,i}|\), where \[Z_{E,i}=\sum_{t=1}^T E(x_t,\widehat s_t)(\widehat s_{t,i}-s_{t,i}).\] The displayed calculation on p. 7 instead identifies this with \(\mathbb E\|Z_E\|_\infty\). These are different. If \(Z\) equals \((1,0)\) and \((0,1)\) with probability \(1/2\) each, then \(\max_i\mathbb E|Z_i|=1/2\), while \(\mathbb E\|Z\|_\infty=1\). This example is realizable with one round, \(S=[0,1]^2\), a fixed zero outcome and a randomized prediction; it refutes the identity, not the final theorem. Also, bounding expected regret to every fixed swap does not alone bound the expectation of the hindsight maximum over swaps.\par
Equation reference: eq:joint. The available proof structure permits a stronger, valid argument. Let \(Q_E=\sum_t E(x_t,\widehat s_t)^2\), \(\Lambda=\log(2d|\mathcal E|T)\) and \(U_T=\sum_t u_t(\widehat s_t,s_t)\). Keeping the realized \(Q_E\) until taking expectations, the outer-regret step gives, simultaneously over coordinates, \[
 |Z_{E,i}|\leq C\bigl(\Lambda+\sqrt{\Lambda Q_E}\bigr)+U_T.\] The residual \(U_T\) is common to every coordinate and event. Taking a maximum over coordinates before expectation is consequently valid on the right-hand side. Jensen’s inequality and the inner-game bound then yield \[\mathbb E\|Z_E\|_\infty
 \leq C\bigl(\Lambda+\sqrt{\Lambda n_T(E)}\bigr)+\mathbb E U_T.\] This argument supplies the missing stronger statement; it does not replace a maximum and expectation by an equality.\par
For a best-response action event \(E_a\), linearity and \(L\)-Lipschitzness give the pathwise inequality \[\mathrm{SwapReg}_T(u)\leq 2L\sum_{a=1}^K\|Z_{E_a}\|_\infty.\] It holds uniformly over all swaps, including the hindsight maximizing swap. Applying Lipschitzness to the event-weighted averages keeps arguments within the convex state space; homogeneity then gives the displayed cumulative expression. Since the action events partition rounds, \(\sum_a n_T(E_a)=T\). Thus, with the corrected inner-loop conventions, the proof yields \(O\bigl(L[K\Lambda+\sqrt{KT\Lambda}]\bigr)\). The factor two is also missing from the original two-term comparison, but has no effect on the asymptotic order.\par
Theorem 4.1’s proof on p. 9 similarly bounds the comparator’s utility error by one coordinate bias, omitting a factor of up to \(d\). With one round, all subsets available, true reward \(\mathbf1\in\mathbb R^d\), prediction zero and event always one, the difference between the two best fixed-action utility totals is \(d\), whereas the coordinate bias is one. A valid pathwise bound for a fixed conditioning function \(\xi\) is \[\mathrm{CReg}_T(\xi,i)\leq
 \sum_{e=1}^d|Z_{E_\xi,e}|+
 \sum_{e=1}^d|Z_{E_{e,\xi},e}|.\] Using [eq:joint] (\#eq:joint) uniformly, with \(Q_E\leq T\), also handles the maximum over the finite conditioning collection. The corrected bound still has order \(\widetilde O(d\sqrt T)\). Thus the local intermediate inequality is false, but its repair preserves the advertised dimension order. These are proof repairs, not empirical tests or counterexamples to the final regret rates with logarithms retained. P2 is Major because the original main decision-making derivations contain invalid expectation and comparator steps. The stronger pathwise argument is a substantive replacement proof, not evidence that the printed proof was valid.\par
\section*{P3: The displayed derivations retain logarithmic factors}
Section 2.2 and Theorem 3.7 drop both the additive logarithm and the square-root logarithm from Theorem 2.4 while using ordinary \(O\), not \(\widetilde O\). The underlying experts theorem indeed contains these terms: the statement of Chen, Luo and Wei’s Theorem 2 was checked in the primary publisher PDF [chen]; its full proof was not independently audited.\par
For multicalibration with \(m\) bins, summing the stated bounds and using Cauchy–Schwarz gives, for a fixed group, \[\mathbb E\mathrm{ECE}(G)
 =O\bigl(T/m+m\Lambda+\sqrt{mT\Lambda}\bigr),
 \qquad \Lambda=\log(2m|\mathcal G|T).\] The printed choice \(m=T^{1/3}\) therefore retains logarithms. Balancing with \(m\) of order \((T/\Lambda)^{1/3}\), using a suitable upper bound for \(\Lambda\), gives \[O\bigl(T^{2/3}\Lambda^{1/3}+T^{1/3}\Lambda^{2/3}\bigr).\] This establishes the exponent up to logarithmic factors, but does not establish the literal log-free headline rate from the supplied derivation. Likewise the swap calculation above retains \(K\Lambda\) and \(\sqrt{KT\Lambda}\). Amend the claims to state the retained factors or supply a sharper argument. No lower bound ruling out log-free rates is asserted here, and the general usefulness of polynomial-time event-unbiased prediction is not negated by this accounting issue. The rounding procedure should also be stated operationally: the adversary must commit before the current private draw is revealed, and the released discretized prediction and its corresponding event calculation need a consistent definition.\par
\section*{Main algorithm checks and smaller corrections}
The copycat inner-game idea is sound under the nonanticipating protocol of p. 4: for a fixed surrogate, \(u_t(s,s)=0\), and the payoff is linear in the adversary’s state. For \(M\) pre-update inner plays \(z_0,\ldots,z_{M-1}\) with respective gradients \(g_0,\ldots,g_{M-1}\), the valid identity is \[\max_{s\in S}\frac1M\sum_{\tau=0}^{M-1}u_t(z_\tau,s)
 =\frac1M\max_{s\in S}\sum_{\tau=0}^{M-1}\langle g_\tau,z_\tau-s\rangle.\] An independent scalar projected-gradient calculation with two threshold events and \(M=10{,}000\) gave both sides \(0.000689125\). It is an illustrative check of the identity, not author-code execution or a general algorithm benchmark.\par
Algorithm 1 samples \(t^2+1\) points including the initial point, but the proof averages \(t^2\) terms and mixes pre- and post-update indices. Make these indices consistent. The additional bounded-payoff point costs only \(O(t^{-2})\) when properly compared, so this is not evidence of a linear-bias failure. Also normalize the outer losses: \(E(\widehat s)(s_i-\widehat s_i)\) can lie in \([-2,2]\), whereas the stated experts primitive assumes \([-1,1]\); dividing by two and restoring the constant suffices. Replace a realized sum of squared events by its expectation only after applying Jensen.\par
The dimension dependence of the inner primitive can be made explicit. At a fixed outer round, put \(a_i=\sum_{\sigma,j}w_{\sigma,i,j}\), so \(\sum_i a_i=1\) and \(|g_i|\leq a_i\). The weighted norm \(\|z\|_a^2=\sum_i a_i z_i^2\) has state-space diameter at most two and dual gradient norm squared at most \(\sum_i g_i^2/a_i\leq1\). Weighted projected gradient descent therefore gives a dimension-independent \(O(\sqrt M)\) regret constant when that projection is available. Zero-weight coordinates can be removed from the norm after projecting the feasible set onto the positive-weight coordinates; a feasible lift must still be supplied for a general state space. On a box, any admissible values in those zero-weight coordinates suffice. With \(M=t^2\), this supplies \(\mathbb E U_T=O(\log T)\), as needed above. An unspecified generic OLO routine can instead have geometry-dependent constants; the implementable projection/optimization access to a general \(S\) should be declared. The simple box state space used for combinatorial rewards supports this construction directly.\par
Two small notation corrections were also verified. Footnote 1 on p. 4 says \(n_T(E)\) is at most the count of rounds where \(E=1\). For \(E\equiv1/2\), \(n_T(E)=T/4\) while that count is zero. The valid statements are \(n_T(E)\leq\mathbb E\sum_t E_t\) and \(n_T(E)\leq\mathbb E\#\{t:E_t>0\}\), with equality to the usual count for binary events. In Theorem C.3 on p. 17 the calibration condition should range over predicted values. As printed over true values, the one-round pair \(s_1=0\), \(\widehat s_1=1\) vacuously satisfies the written equality, but fails transparency for \(u(s)=s\). The surrounding term “fully calibrated” and the proof make the intended correction clear. Under the corrected calibration condition and the standing linear-utility assumption, the partition/coarsening arguments of C.3–C.8 are consistent.\par
\section*{Information timing, repair scope and limitations}
Section 2 explicitly lets the adversary see the prediction distribution before committing the state, and only then samples the realized prediction. This restriction is essential. Section 4 lists action commitment before reward production; specify that the current private action or realized public prediction is not revealed to the adversary before that commitment. If an adversary could see a realized singleton action, it could assign that action reward zero and the other reward one. The learner would get zero while the best fixed action gets at least \(T/2\), precluding the claimed no-regret guarantee. Exhaustive enumeration of all \(2^8\) action sequences verifies this diagnostic for \(T=8\). This stronger observation model is excluded by Section 2, so it is an interpretation warning, not a counterexample under the intended model.\par
Appendix D additionally needs an optimization-access and representation statement, including treatment of lower-dimensional feasible sets and precision. Bounded coefficient magnitude alone does not bound rational encoding length; an explicit rounding-error budget is required. These computational details were not independently completed. The null-event counterexample and core proof defects above are established; a fully repaired accelerated implementation and its general oracle runtime remain evidence requests.\par
All principal displayed proofs and relevant appendix algorithms were read. Source pp. 3–9 and 17–22 were actually viewed at legible resolution; the paper contains no source plots. Appendix A’s related-work discussion and the bibliography were not verified for priority or completeness. No source repository, executable implementation, experimental data or historical run was available to inspect. Author adoption is unknown. Independent checking is recorded in the verification files. The demonstrated corrections preserve the main reductions up to logarithmic factors under the stated primitive and timing assumptions; exact-rate and accelerated-method claims require the identified amendments.\par
2 G. Noarov, R. Ramalingam, A. Roth and S. Xie. High-Dimensional Prediction for Sequential Decision Making. PMLR 267 (2025), 46762–46783. https://proceedings.mlr.press/v267/noarov25b.html. L. Chen, H. Luo and C.-Y. Wei. Impossible Tuning Made Possible: A New Expert Algorithm and Its Applications. PMLR 134 (2021), 1216–1259, Theorem 2. https://proceedings.mlr.press/v134/chen21f/chen21f.pdf.\par
\section*{Scope and limitations}
Author programs were not executed. Scientific adjudication remains pending; human reviews are empty.\par
\end{document}
